% Generated draft: requires mathematical and source review.
\documentclass[11pt]{article}
\usepackage[margin=25mm]{geometry}
\usepackage{fontspec}
\setmainfont{Latin Modern Roman}
\usepackage{amsmath,amssymb,mathtools,unicode-math}
\setmathfont{Latin Modern Math}
\usepackage{enumitem,needspace,xurl,hyperref}
\setlength{\parindent}{0pt}
\setlength{\parskip}{5pt}
\setlength{\emergencystretch}{3em}
\newcommand{\N}{\mathbb{N}}
\newcommand{\Z}{\mathbb{Z}}
\newcommand{\Q}{\mathbb{Q}}
\newcommand{\R}{\mathbb{R}}
\newcommand{\C}{\mathbb{C}}
\newcommand{\F}{\mathbb{F}}
\newcommand{\A}{\mathbb{A}}
\newcommand{\PP}{\mathbb{P}}
\newcommand{\E}{\mathbb{E}}
\newcommand{\eps}{\varepsilon}
\newcommand{\dd}{\,\mathrm{d}}
\begin{document}
\section*{UnsolvedMath ID: 30000627}
\textbf{Source problem code:} catalogue entry
\subsection{Definitions and mathematical statement}
Let \(K\subset\R^3\) be a compact convex body with nonempty interior and sufficiently differentiable boundary. At \(x\in\partial K\), use the nonnegative convex-surface sign convention for principal curvatures, ordered \(\kappa_{\min}(x)\leq\kappa_{\max}(x)\). Define the ordered principal curvature radii by
\[
R_1(x)=\kappa_{\max}(x)^{-1},\qquad
R_2(x)=\kappa_{\min}(x)^{-1},
\]
with reciprocal zero interpreted as infinity if zero curvature is permitted. The source's Alexandrov uniqueness assertion is
\[
\exists C\in(0,\infty)\quad\forall x\in\partial K,\quad
R_1(x)\leq C\leq R_2(x)
\quad\Longrightarrow\quad K\text{ is a Euclidean ball}.
\]
The constant is the same at every boundary point. Equivalently, with the stated extended-radius convention, the hypothesis is
\[
\kappa_{\min}(x)\leq C^{-1}\leq\kappa_{\max}(x)
\qquad(x\in\partial K).
\]
If the conclusion is a ball, its radius must equal \(C\), so the precise conclusion can be written \(K=a+C\overline{B^3}\) for some \(a\in\R^3\).

Smoothness is not quantified in the record. Pointwise curvatures require at least a suitable twice-differentiable setting, while \(C^2\), \(C^\infty\), and real-analytic versions of this rigidity assertion must not be treated as the same hypothesis. The draft does not add analyticity or positive Gaussian curvature. If the source instead requires all radii finite, that imposes strict positivity of both principal curvatures and must be stated separately; the literal convex-body wording does not explicitly impose it.

The supplied literature triage checked 2026-08-21 labels the extracted assertion open and reports no authoritative complete resolution located, but gives no separate regularity-specific theorem or reference. That generic assessment is not independent certification that the literal smooth formulation is unresolved or correctly stated. Source review should check the exact regularity and any omitted hypotheses before presenting this classical named assertion as a current conjecture.
\subsection{Short English statement}
The source asserts that a convex surface whose two curvature radii straddle one fixed positive number at every point must be spherical. Its exact smoothness and zero-curvature conventions are essential and are not fixed by the supplied open-status label.
\subsection{Sources}
\begin{itemize}
\item Local UnsolvedMath registry, permanent ID 30000627, catalogue entry, \emph{Alexandrov Uniqueness Conjecture for Convex Surfaces}; attribution to Alexandrov is explicit in the original wording.
\item \emph{Konvexgeometrie} (2007), Oberwolfach Reports, pp. 3321--3396, DOI \url{https://doi.org/10.4171/owr/2006/56}.
\item Supplied literature triage checked 2026-08-21 cites only that report; no regularity-specific proof, disproof, or complete theorem statement is supplied.
\end{itemize}
\end{document}
